\bookchapter{Common Action, Moving Atoms and Fully Paid Molecular Binding}

\subsection{Particle motion and field energy on the same geometry}

Occupied electronic states connect to Volume I through a common action. The completed recipient places the atom and its constrained field on one explicit radial geometry. Write the constant lapse as $\mathcal N$, with $0<\mathcal N\leq1$, and the spatial metric and volume measure as
\[
\begin{gathered}
ds^2=\mathcal N^{-2}dr^2+r^2d\Omega^2,
 \\
 d\mathrm{vol}=\mathcal N^{-1}r^2dr\,d\Omega.
\end{gathered}
\]
The areal sphere remains $4\pi r^2$. A lapse-weighted particle action with action unit $I_*$ gives the external kinetic form
\[
 \tau\mathcal N^3|\partial_r\psi|^2
       +\frac{\tau\mathcal N}{r^2}|\nabla_\Omega\psi|^2.
\]
Both factors follow from this metric and measure. A fixed-position clock multiplying an unchanged Hamiltonian by $\mathcal N$ is another specified recipient; the present atom includes the changed spatial geometry. \samref{radial_geometry}

The constrained scalar and polarization field uses that same geometry. Its positive Green inverse has angular exponent
\[
 \alpha_l=\frac{-1+\sqrt{1+4l(l+1)/\mathcal N^2}}2,
\]
\[
 G_{\mathcal N}(x,y)=\frac1{2\pi}
  \sum_{l\geq0}\frac{2l+1}{2\alpha_l+1}
  \frac{r_<^{\alpha_l}}{r_>^{\alpha_l+1}}P_l(n_x\cdot n_y).
\]
At $\mathcal N=1$ this is the original Coulomb Green kernel. Its monopole is unchanged, so the uniform nuclear shell retains its core self and scalar potential. Higher angular terms, pair interaction and exchange use the new complete kernel. The all-degree remainder is explicitly bounded, and the finite field determinant stays in the reference.

Three geometries, $\mathcal N=1,3/4,1/2$, produce $132$ positive populated binding accounts on the original helium/carbon sources. Every account has an occupied witness. Thirty numerical ionization lower bounds are positive; $102$ intervals contain zero. The complete source phase and self terms remain inside both nuclear minima. This geometry makes the common radial units explicit:
\[
\begin{gathered}
\ell_{\rm rad}=\frac{\tau\mathcal N^3}{v_c},\\
 E_{\rm rad}=\frac{v_c^2}{\tau\mathcal N^3},\\
 t_*=I_*/\epsilon_*.
\end{gathered}
\]
A physical lapse, action and absolute field radius are the concrete remaining inputs to their physical assignment.

\subsection{Full massive action and collective recoil}

The local massive action identifies $\tau=I_*^2c_*^2/(2E_0)$. With full geometric Laplacian $L_g$, its rest-subtracted external kinetic operator is
\[
 T_{\rm rel}=\mathcal N E_0
 \left[\sqrt{1+2\tau L_g/E_0}-1\right].
\]
Its low-momentum quadratic recipient has exactly the preceding radial and angular factors. The square root acts on the full spatial space before an orbital matrix is formed. This ordering carries momentum states outside the retained orbital span into the occupied source. \samref{common_action}

The completed common-action panel evaluates the $132$ occupied accounts at six rest regimes, giving $792$ positive total-binding accounts. Every occupied upper witness improves. The full fourth moment pays smooth inner states; compact hard-support shells use their entire momentum tail, including their endpoints. For every retained occupied account, $E_0/E_{\rm atom}\geq1167221$ is a sufficient regime for a complete kinetic correction of at most one percent of that state's nonrelativistic kinetic energy. This is a calculated finite-family sufficient regime. The six evaluated rest ratios $8,32,128,512,8192,131072$ retain their own larger or smaller enclosed errors.

Finite nuclear mass adds collective electron recoil. Let $\Pi$ project onto the actual one-electron orbital span, $P_\alpha=\Pi p_\alpha\Pi$ and $Q=\Pi|p|^2\Pi$. The positive remainder
\[
 R=Q-\sum_\alpha P_\alpha^2
   =\sum_\alpha\Pi p_\alpha(1-\Pi)p_\alpha\Pi
\]
records momentum intermediates outside that span. On the full fermion sector, with $\Gamma$ denoting second quantization, the literal compressed collective second moment is
\[
 M_2(P)=\sum_\alpha(P_\alpha^{\rm total}I-\Gamma(P_\alpha))^2+\Gamma(R).
\]
It retains cross-electron momentum and exchange. Total atomic momentum zero still leaves internal core recoil. \samref{recoil_moment}

The full fourth-moment matrix $Q_4$ yields $M_4(0)\leq B_4=N_e^3\Gamma(Q_4)$. For the massive core kinetic $K_m=\sqrt{m^2+a^2p^2}-m$, spectral calculus gives
\[
 \frac{a^2M_2}{2m}-\frac{a^4B_4}{8m^3}
 \leq\Pi K_m\Pi\leq\frac{a^2M_2}{2m}.
\]
This encloses the true square root before compression. The completed scale-one atom retains all $256$ configurations, all $N_e=0,\ldots,8$ sectors and both original Gray phases. The largest midpoint error is $2.861\times10^{-7}$ source energy units. The corresponding twenty-seven contact controls have instrument-error upper bound $7.250\times10^{-8}$. H002381 also corrects the earlier numerical scale association: the frozen H002378 matrices used the scale-two orbital family with unscaled coefficients; their matrix results remain, while the literal scale-one recipient supplies the corrected same-atom association.

\subsection{History acts on the full occupied atom}

A retained isotope history supplies actual phases and contact times. Between contacts the full moving Pauli Hamiltonian $H_s$ acts on occupied states, ion sectors and free continua. For an original Gray change $s\to s'$, the complete difference is bounded:
\[
\begin{gathered}
H_{s'}-H_s=\Delta A=\delta\chi(p_c)-V,
 \\
 0\leq\chi\leq1,
 \\
 \|V\|\leq\frac{N_e\sqrt2}{10t}.
\end{gathered}
\]
Kinetic and pair terms remain in each full Hamiltonian even when identical entries cancel in the difference. Common domains give complete propagators. For history $h$ with phases $s_0,\ldots,s_n$ and intervals $\Delta t_i$, the conditional atomic evolution is
\[
 U_h=e^{-iH_{s_n}\Delta t_n/I_*}\cdots
       e^{-iH_{s_0}\Delta t_0/I_*}.
\]
The actual order matters. Averaging $U_h\rho U_h^*$ over the retained history law gives a completely positive trace-preserving driving channel; correlated preparations retain their conditional density in that average. \samref{history_atomic}

The full energy change is the sum of incoming-state switch work, $w_i=\operatorname{Tr}(\rho_i\Delta A_i)$. Because the source contact already pays its nuclear gap $\delta_i$, the extra atomic payment is $w_i-\delta_i$. The recoil allowance bounds this extra term using the actual incoming recoil moment. This prevents counting the same source energy twice.

The coupled quantum contact supplies the corresponding full atomic-plus-field instrument. For the two eligible phases put $\Gamma_s=e^{-\beta(A_s+W_s)}$ and $\Gamma=\Gamma_a+\Gamma_b$. Bounded operators satisfy
\[
\begin{gathered}
K_s\Gamma^{1/2}=\Gamma_s^{1/2},\\
 \sum_sK_s^*K_s=I,\\
 K_s\Gamma K_s^*=\Gamma_s.
\end{gathered}
\]
Both changed and held outcomes alter the conditional atomic state. The paired channel retains thermal weight and KMS reciprocity; its scalar limit is the original complete-energy heat-bath rule. Its complete mean energy increase obeys
\[
 \sum_s\operatorname{Tr}[K_s\rho K_s^*(A_s+W_s)]
 \leq\operatorname{Tr}[\rho(A_a+W_a)]+\frac{\log2}{\beta}.
\]
The full moving-space result uses thermal weights on their proper domains, including the free continuum. \samref{quantum_contact}

An explicit carrier receives the opposite energy change and retains the outcome record. Its energy intertwining uses the same $I_*$. A finite normal carrier packet has a quantified energy-class accuracy bound, with its width and energy variance both charged. In the retained Gaussian example, width $10^{-6}$ source time gives a trace-norm bound below $0.025598$ and energy variance $2.5\times10^{11}$ source units squared. These are completed preparation operands; the positive-energy field realization and its original correlations have their separate receiving construction.

\subsection{An actual eight-electron original contact preparation}

The original-frame calculation joins that contact to a literal atom. It uses frame $0$ to frame $7$, mark $6$, all forty-eight signs, $U_a=(1,1)$, $U_b=(0,2)$, $C_a=25361/10$, $C_b=11391/5$, and $\beta=1/16$. The electron preparation is the original four spatial orbitals with both spins occupied, retaining all $264$ angular/radial components. Its axis is along $U_b$. The total-momentum preparation is the normalized Gaussian
\[
\begin{gathered}
\Omega(P)=\pi^{-3/4}e^{-|P|^2/2},\\
 \langle P\rangle=0,\quad\langle P^2\rangle=3/2.
\end{gathered}
\]
Full massive electron kinetic, collective recoil and the original dressed $T$ are retained. The preparation consequently has finite localization and motion. \samref{actual_atom}

The detached reference has the same Gaussian, source record, held field and baseline. Full source recoil and the attached-minus-detached recoil difference have distinct places in the ledger. The actual intrinsic energies are $2504$ and $2246$; conditional-minimum orbital-bank rows both using $440$ cannot substitute for these records. Record-independent native integrals are reused and the actual phase-dependent mass and orientation are compiled from them.

A complete-space lower gauge retains positive pair repulsion. With $x_i=1/\max(|r_i|,r_c)$, the pair bound $|r_i-r_j|^{-1}\geq x_ix_j/(x_i+x_j)$ gives a concave box function. Its vertex with $k$ nonzero coordinates has energy
\[
 f_k=\frac{-(41/5)k+(9/40)k(k-1)}{r_c}.
\]
For $N_e=8,r_c=1$, its minimum is $-53$. With original wave offset $-11/10$, the common floor is $-541/10$. Actual electronic means are approximately $-7.9399362574$ and $-8.1451906120$; full Gaussian/electron recoil is below $0.000231495$. The outgoing source energy above the common gauge is at most $44.855040883$.

Spectral Jensen and the full source-normalized cube-root Gibbs receiver yield selected-contact acceptance at least $0.04714669$ across all sixteen original directions. This is a full joint bound on the supplied eight-electron preparation. The full joint mean/effect matrix and both ordered mixed covariances are the named next recipient.

The completed original-boundary supplement also supplies a finite-energy transverse datum. A field displacement reaches $\xi=0$ while preserving the full source readouts; its minimum position preparation cost is $851358848/7425$. A momentum field of kinetic energy one gives $\dot\xi=\sqrt{2V}>0$. The common preparation work is charged once in source and matched reference and in their common gauge, so the energy-above-gauge and acceptance bounds are unchanged. This makes explicit why an eligible original boundary and a paid atomic preparation belong in the same ledger.

\subsection{One molecular field and the true dissociation reference}

Two atoms must inhabit one field geometry. Put their centers at $R_A=-De/2$, $R_B=De/2$ and define $r_a=|x-R_a|$, $n_a=(x-R_a)/r_a$. The completed two-center recipient is
\[
 n_D(x)=\frac{r_Bn_A+r_An_B}{|r_Bn_A+r_An_B|}.
\]
It is translation/rotation covariant, exchange symmetric and of unit norm. Near either core it tends to the original radial recipient at large separation. Unit norm keeps the electron's regularized polarization self constant through assembly. The source's exact Cartesian $U$ remains necessary in cross terms; $|U|^2$ alone does not encode orientation. \samref{molecular_geometry}

The molecular Hamiltonian includes both original intrinsic energies, both finite core self energies, full electron massive kinetic, every scalar/polarization core-electron attraction, every pair/exchange term and the core cross energy
\[
 C_{AB}(D)=\frac{g_c^2Z_AZ_B/v_c+g_p^2U_A\cdot U_B/v_p}{D}.
\]
Changing $D$ also changes $n_D$. Its derivative therefore acts in the own-core potential and internal electron-pair terms. Those terms contribute at the same leading order as cross attraction. Their full signed response is retained before its sign is optimized.

Antisymmetric products of compactly localized actual atomic states give explicit admissible molecular trials. Disjoint supports eliminate cross Coulomb exchange while preserving all internal correlation and exchange. The nonlocal massive one-body kinetic still adds through the block-diagonal one-particle density. Localization error and recipient Taylor remainder are paid. The separated ground is minimized over every electron allocation, every free-electron channel, every original nuclear phase and all fragment momenta. A prepared neutral pair and that complete dissociation reference are distinct quantities, both retained in the source record.

\subsection{Complete moving-reference attraction for carbon and oxygen}

For carbon the complete moving comparison retains ninety-one charge/free allocations and all fragment momenta. Its phase condition is $\lambda r_c>923/2500$. In this region minimizing nonvacuum atoms are attained below their escape thresholds, and no minimizing allocation leaves free electrons. All positive finite core rest parameters $B$ are admitted. The earlier sufficient heavy-core condition remains its quantitative predecessor; the completed signed-source argument removes it from this existence result. \samref{carbon_moving}

For a correlated ion write
\[
\begin{gathered}
v=U-\mu,\\
 t=TU-V_{\rm int},
\end{gathered}
\]
where $\mu$ is the complete electron direction mean, $T$ the finite-core own-field response and $V_{\rm int}$ the changing internal polarization-pair response. Carbon's signed leading coefficient is
\[
 C_1(e)=(6-n_A)(6-n_B)+v_A\cdot v_B/10
                         +e\cdot(t_A-t_B)/10.
\]
Core/state reversal at the same fragment momentum preserves atomic energy and reverses the appropriate signed moments. Optimizing signs and separation gives $C_1\leq-(n_A-6)^2$ for unequal charges and $C_1\leq-18/5$ for matched neutral atoms. Thus $C_1\leq-1$ suffices over the complete minimizing reference.

The actual four-helium $8p8n$ source has a parallel completed result with all $2^{48}$ phases, $83$ receiver classes and $153$ charge/free allocations. Its phase condition is $\lambda r_c>2/15$. All twenty-eight higher-$q$ competitors have positive paid margins already at that boundary; the smallest is about $9.0812315001$. Its leading coefficient replaces $6$ by $8$ above, giving $C_1\leq-1$ across all seventeen no-free charge splits and $-32/5$ for matched neutral conjugate states. Possible overcharged minimizing ions remain in the complete reference. \samref{oxygen_moving}

Writing $C_*=344\lambda+224/(5r_c)$ and $\tau_c=\tau E_0/(B+C_*)$, oxygen's full molecular trial obeys
\[
 \begin{aligned}
 E_{\rm mol}(0)\leq{}& S_B(0)-\frac1{D+L}\\
 &+\frac{3\pi^2(16\tau+256\tau_c)}{2R^2}
 +\frac{22\tau_c}{L^2}\\
 &+\frac{721000000R+80L}{(D-L)^2}.
 \end{aligned}
\]
The terms pay actual-ground preparation, collective localization, relative recoil and every common-field/direction remainder. Choosing $D=k^3$, $L=k^2$, $R=k^2/128$ gives positive gain exceeding $1/(2k^3)$ for sufficiently large finite $k$ at every fixed positive finite $B$. Actual ion-gap magnitudes and resulting bond separation remain unevaluated. The completed result is strict molecular lowering below the complete moving reference at finite separation.

\subsection{Feedback preparation and spatial formation are paid}

The nuclear contact also has a positive-stiffness Gaussian field realization. For $m$ cores, $U=\sum_i u_i$, $q_i=|u_i|^2$ and $J=g^2/\kappa$, its augmented energy is
\[
 H_{\rm aug}=\sum_iE_i+\frac J2\sum_iq_i
 +\frac\kappa2|h|^2-gh\cdot U
 =I+\frac\kappa2|h-(g/\kappa)U|^2.
\]
The normalized Gaussian trace recovers the original contact energy $I$ exactly. Stored energy, drive and diagonal counterterm all stay. Within the declared Gaussian class, the collective two-component field minimizes independent-core preparation payment. Its attaining Gram has rank two; edge-by-edge realizations cost twice as much for three cores and three times as much for four. \samref{paid_feedback}

The same source identifies $\epsilon_{\rm phase}=G_{\rm fb}/R$ and $\lambda r_c=G_{\rm fb}/G_{\rm scalar}$. The complete spatial envelope gives favorable carbon formation for $d<13R$, with barrier $2/3$ in $G_{\rm scalar}/R$, and favorable four-core formation for $d<19.5R$, with barrier $8/9$. All $2^{36}$ and $2^{48}$ phases and their changing ground multiplicities remain. The original ratio $G_{\rm fb}/G_{\rm scalar}=1$ satisfies a common sufficient condition greater than $2/5$ for both bare formation and moving molecular regions.

Field protocols include their actual preparation. A separated source whose initial field already carries each core has another payment from an initially free collective field, even if their effective phase dynamics agree. At infinite separation the former field KL is $\beta m\sigma$, twice the independent-core free-field payment. This difference is determined by the specified initial state. The completed twenty-four protocols retain full work, free energy and field relative entropy. Their finite-temperature free-formation signs change with temperature while every protocol has its positive relaxation payment.

\subsection{A normalized virtual packet with every molecular motion retained}

The full molecular-motion continuation goes beyond clamped core coordinates. It keeps the original autonomous Hamiltonian, original electronic preparation and all ninety-one carbon fragment channels. Let $G$ contain every matched opposite source minimum, $P=G\otimes|0\rangle$ its vacuum sector, and $S$ the complementary paired source space with the full Fock space. For any normalized constant vector $c$ in the entire ground space, define
\[
\begin{gathered}
R_S=(H_0|_S)^{-1},\\
 u_c=tR_SSA Pc,\\
 F_c(x)=\frac{Pc+u_c(x)}{\sqrt{1+\|u_c(x)\|^2}},\\
 t=q(R)/180.
\end{gathered}
\]
The molecular trial multiplies this state by the original cosine packets and electronic state and applies the full source-controlled Weyl displacement. It retains $12{,}259$ paired ground states in thirty-six complete factors and all $910$ departures. The virtual tail contains the complete source-conditioned Fock resolvent, oscillator counts, STF space and Gaussian continuum frequencies. \samref{molecular_motion}\samref{motion_derivation}

Full compound-Poisson inverse returns obey
\[
 \frac1{\Delta+\mu}\leq R_{\rm self}
 \leq\frac1{\Delta+\mu}
       +\frac{\operatorname{Var}}{\Delta(\Delta+\mu)^2}.
\]
Reciprocal entries remain in the complete return operator $B(x)$, with uniform $B_-I\leq B(x)\leq B_+I$ across the original packet support. The source gap and full field moments pay the virtual norm and its coordinate derivatives, using
\[
 \partial_xR_S=-R_S(\partial_xH_S)R_S.
\]
The entire resolvent is differentiated, including source offsets, frequency changes and Weyl profiles.

The ground packet and virtual tail have orthogonal source labels. All original kinetics and dressing connections are source-diagonal, so their cross kinetic terms vanish exactly. Pointwise normalization preserves the common electronic terms and supplies its exact derivative metric. Radial/core connection squares, full field momentum, original electron recoil and the original massive core actions are then paid. Tail-tail exchange and normalization remain in
\[
\begin{gathered}
G_{\rm source}=\frac{t_-^2B_--t_+a n_*}{1+n_*},
 \\
 G_{\rm motion}=G_{\rm source}-\Delta T_{\rm radial}-\Delta T_{\rm core}.
\end{gathered}
\]
All six original carbon families retain positive full-motion gains. For 333, 343 and 444 at weak polarization they are approximately $7.19359,7.75886,8.90615$ times $10^{-6}$ source units; at strong polarization they are $7.31984,8.99286,12.37149$ times $10^{-6}$. The largest added motion payment is $7.02436\times10^{-7}$. Radial motion dominates that increment; massive core and shared-field recoil increments are at most $3.128\times10^{-12}$.

All thirty-six central molecular witnesses improve; the least new full-fragment binding lower bound is $356.88838984$ source units. The $1548$ original preparation rows receive the additional ground witness for their own family and rest chart. Their historical prepared energies remain unchanged. All thirty-six complete references, or $3276$ explicit fragment-channel occurrences, are retained. Three native scientific returns, $40{,}327$ independent assertions, $150$ infinite Laplace quadratures and six derivative checks close this full-motion construction.

The finished result is a fully paid molecular ground-state variational gain, classified strong contact. The clamped exact-eigenvalue intervals retain their own completed meaning. Its physical receiving equations are $\Delta E_{\rm phys}=E_*G_{\rm motion}$, $t_{\rm phys}=(I_*/E_*)t_{\rm model}$ and $x_{\rm phys}=L_*x_{\rm model}$. The same constituent/fragment reference must supply $E_*$, the source/current phase supplies $I_*$, and the spatial source/receiver supplies $L_*$. These are the precise remaining common physical assignments joining atomic attachment, nuclear formation, molecular motion and measured energy, time and length.
